% !TEX program = xelatex
% Paper Pilot - 中国研究生学位论文模板
% 编译方式: xelatex main.tex -> bibtex main -> xelatex main.tex -> xelatex main.tex

\documentclass[12pt,a4paper,openany]{ctexart}

% ── 页面设置 ──
\usepackage[top=2.54cm, bottom=2.54cm, left=3.17cm, right=3.17cm]{geometry}
\usepackage{setspace}
\onehalfspacing

% ── 数学 & 表格 ──
\usepackage{amsmath,amssymb,amsfonts}
\usepackage{booktabs}
\usepackage{tabularx}
\usepackage{multirow}
\usepackage{graphicx}
\usepackage{float}

% ── 代码 ──
\usepackage{listings}
\lstset{
  basicstyle=\small\ttfamily,
  breaklines=true,
  frame=single,
  numbers=left,
  numberstyle=\tiny\color{gray},
}

% ── 超链接 ──
\usepackage{hyperref}
\hypersetup{
  colorlinks=true,
  linkcolor=black,
  citecolor=blue,
  urlcolor=blue,
  bookmarksnumbered=true,
}

% ── 页眉页脚 ──
\usepackage{fancyhdr}
\pagestyle{fancy}
\fancyhf{}
\fancyhead[C]{\small\leftmark}
\fancyfoot[C]{\thepage}
\renewcommand{\headrulewidth}{0.4pt}

% ── 参考文献样式 (GB/T 7714) ──
% 使用时取消下行注释，需安装 gbt7714-bibtex-style
% \bibliographystyle{gbt7714-numerical}

% ── 信息设置 ──
\newcommand{\thesisTitle}{{{TITLE}}}
\newcommand{\thesisAuthor}{{{AUTHOR}}}
\newcommand{\thesisDegree}{{{DEGREE}}学位论文}
\newcommand{\thesisDate}{\today}
\newcommand{\thesisUniversity}{XX大学}

\begin{document}

% ════════════════════════════════════════
% 封面
% ════════════════════════════════════════
\begin{titlepage}
\centering
\vspace*{2cm}

{\Huge\bfseries \thesisUniversity}

\vspace{2cm}

{\LARGE\bfseries \thesisTitle}

\vspace{3cm}

{\Large
\begin{tabular}{rl}
学\qquad 院： & XX学院 \\[0.5em]
专\qquad 业： & XX专业 \\[0.5em]
学\qquad 生： & \thesisAuthor \\[0.5em]
指导教师： & XX教授 \\[0.5em]
\end{tabular}
}

\vfill

{\large \thesisDate}

\vspace{1cm}
\end{titlepage}

% ════════════════════════════════════════
% 中文摘要
% ════════════════════════════════════════
\newpage
\begin{center}
{\Large\bfseries 摘\quad 要}
\end{center}
\vspace{0.5em}

\noindent\textbf{摘要：}这里是中文摘要的内容。请在此处撰写论文的研究背景、研究方法、主要结果和结论。摘要应简明扼要，一般不超过500字。

\vspace{1em}
\noindent\textbf{关键词：}关键词1；关键词2；关键词3；关键词4

\vspace{2em}
\hrule
\vspace{1em}

\begin{center}
{\Large\bfseries Abstract}
\end{center}
\vspace{0.5em}

\noindent\textbf{Abstract:} Write your English abstract here. Briefly describe the research background, methodology, main results, and conclusions.

\vspace{1em}
\noindent\textbf{Keywords:} keyword1; keyword2; keyword3; keyword4

% ════════════════════════════════════════
% 目录
% ════════════════════════════════════════
\newpage
\tableofcontents
\newpage
\listoffigures
\newpage
\listoftables

% ════════════════════════════════════════
% 正文
% ════════════════════════════════════════
% 第一章 绪论
\section{研究背景与意义}

\subsection{研究背景}
\label{sec:background}

在此处描述研究领域的背景信息。随着XXX技术的发展（如公式~\ref{eq:example}），相关问题日益突出。

\subsection{研究现状}
\label{sec:current}

国内外学者在XXX领域已有大量研究~\cite{zhang2023,l wang2022}。图~\ref{fig:example}展示了该领域的发展趋势。

\begin{figure}[H]
  \centering
  % \includegraphics[width=0.8\textwidth]{figures/example.pdf}
  \fbox{\parbox{0.8\textwidth}{\centering\vspace{2cm}图片占位符\vspace{2cm}}}
  \caption{示例图片}
  \label{fig:example}
\end{figure}

\subsection{研究目标与内容}
\label{sec:objectives}

本文的主要研究目标包括：
\begin{enumerate}
  \item 目标一：XXX
  \item 目标二：XXX
  \item 目标三：XXX
\end{enumerate}

表~\ref{tab:example}展示了本文的研究内容概览。

\begin{table}[H]
  \centering
  \caption{研究内容概览}
  \label{tab:example}
  \begin{tabular}{ccc}
    \toprule
    章节 & 研究内容 & 方法 \\
    \midrule
    第二章 & 文献综述 & 综合分析法 \\
    第三章 & 方法设计 & 实验法 \\
    第四章 & 结果分析 & 统计分析法 \\
    \bottomrule
  \end{tabular}
\end{table}

\subsection{论文组织结构}
\label{sec:structure}

本文共分为五章，各章内容安排如下：
\begin{itemize}
  \item 第一章：绪论，介绍研究背景、现状与目标
  \item 第二章：文献综述，梳理相关研究
  \item 第三章：研究方法
  \item 第四章：实验结果与分析
  \item 第五章：总结与展望
\end{itemize}

示例公式如下：
\begin{equation}
  E = mc^2
  \label{eq:example}
\end{equation}

% 第二章 文献综述
\section{相关理论基础}
\label{sec:theory}

\subsection{XXX理论}
\label{sec:theory_xxx}

在此处描述相关理论基础。

\subsection{YYY方法}
\label{sec:method_yyy}

在此处描述相关研究方法。

\section{国内外研究进展}
\label{sec:literature}

\subsection{国外研究进展}
\label{sec:foreign}

引用国外文献~\cite{smith2021}。

\subsection{国内研究进展}
\label{sec:domestic}

引用国内文献~\cite{zhang2023}。

\section{本章小结}
\label{sec:lit_summary}

本章梳理了XXX领域的理论基础与研究进展，为后续研究奠定了基础。

% 第三章 研究方法
\section{研究设计}
\label{sec:design}

\subsection{实验方案}
\label{sec:experiment}

\subsection{数据采集}
\label{sec:data}

\section{关键技术}
\label{sec:techniques}

\section{本章小结}
\label{sec:method_summary}

% 第四章 实验结果与分析
\section{实验结果}
\label{sec:results}

\section{结果分析}
\label{sec:analysis}

\section{讨论}
\label{sec:discussion}

\section{本章小结}
\label{sec:res_summary}

% 第五章 总结与展望
\section{研究总结}
\label{sec:conclusion}

本文的主要工作和创新点包括：

\begin{enumerate}
  \item 创新点一：XXX
  \item 创新点二：XXX
  \item 创新点三：XXX
\end{enumerate}

\section{研究不足}
\label{sec:limitations}

\section{未来展望}
\label{sec:future}


% ════════════════════════════════════════
% 参考文献
% ════════════════════════════════════════
\newpage
\bibliographystyle{plain}
\bibliography{references}

% ════════════════════════════════════════
% 附录
% ════════════════════════════════════════
\newpage
\appendix
\chapter{附录A：补充数据}
这里是附录内容。

% ════════════════════════════════════════
% 致谢
% ════════════════════════════════════════
\newpage
\begin{center}
{\Large\bfseries 致\quad 谢}
\end{center}
\vspace{1em}

在此感谢导师的悉心指导，感谢同学的帮助与支持。

\end{document}
